\documentclass[11pt]{article}
\usepackage[utf8]{inputenc}
\usepackage{amsmath,amssymb}
\usepackage[margin=1in]{geometry}
\usepackage{hyperref}
\emergencystretch=3em

\title{Every term at $d=3$:\\
$Z_3(n)=n^{-1/2}\,P(\log n)+O(n^{-1})$ with $P$ an explicit quadratic}
\author{Claude, with Daniel Murfet}
\date{2026-09-06}

\begin{document}
\maketitle
\sloppy

\begin{abstract}
Unit~38 gave the leading term $Z_3(n)\sim\tfrac A8n^{-1/2}(\log n)^2$ for the $d=3$ chart integral
with three coinciding exponents. Iterating the tail mechanism of unit~41 once more identifies the whole
polynomial: with $\theta(x)=H(x)-A\log x-C_0\in[0,M/x]$ for $x\ge1$,
\[
H_2(L)=\tfrac A2\log^2L+C_0\log L+C_1-\int_L^\infty\frac{\theta(x)}{x}\,dx,\qquad
0\le\int_L^\infty\frac{\theta}{x}\le\frac ML,
\]
with $C_1=\int_0^1H/x+\int_1^\infty\theta/x$, and hence $Z_3(n)=n^{-1/2}P(\log n)+O(n^{-1})$ where
$P(\ell)=\tfrac A8\ell^2+(\tfrac32A\log b+\tfrac12C_0)\ell+(\tfrac92A\log^2b+3C_0\log b+C_1)$.
Zero \texttt{sorry}/\texttt{axiom}.
\end{abstract}

\section{Setting}
The recursion $H_{m+1}(L)=\int_0^LH_m(x)/x\,dx$ maps an expansion ``polynomial in $\log x$ plus a
remainder $O(1/x)$'' to an expansion of the same shape one degree higher: the polynomial part
integrates exactly against $dx/x$ over $[1,L]$, the remainder over $[1,\infty)$ converges to a new
constant, and its tail over $[L,\infty)$ is again $O(1/L)$. This unit runs the step once, from $H$
(unit~41) to $H_2$.

\section{The things formalised}
{\small
\begin{verbatim}
noncomputable def logRemainder (β a x : ℝ) : ℝ :=
  logPrimitive β a x - quadMass β a * Real.log x - logConst β a
theorem logRemainder_div_integrableOn ... (hc : 1 ≤ c) :
    IntegrableOn (fun x => logRemainder β a x / x) (Ioi c)
noncomputable def logConst2 (β a : ℝ) : ℝ :=
  (∫ x in Ioc 0 1, logPrimitive β a x / x) + ∫ x in Ioi 1, logRemainder β a x / x
theorem logLogPrimitive_sub_poly ... (hL : 1 ≤ L) :
    logLogPrimitive β a L
      - (quadMass β a * Real.log L ^ 2 / 2 + logConst β a * Real.log L + logConst2 β a)
      = -∫ x in Ioi L, logRemainder β a x / x
theorem threeDIntegral_second_order_isBigO ... :
    (fun n => threeDIntegral β a b n - n ^ (-(1/2 : ℝ)) * threeDSecondOrderCoeff β a b n)
      =O[atTop] fun n => n⁻¹
theorem threeDIntegral_second_order_tendsto ... :
    Tendsto (fun n => Real.sqrt n * threeDIntegral β a b n
        - quadMass β a / 8 * Real.log n ^ 2
        - (3 * quadMass β a * Real.log b / 2 + logConst β a / 2) * Real.log n)
      atTop (𝓝 (9 * quadMass β a * Real.log b ^ 2 / 2 + 3 * logConst β a * Real.log b
        + logConst2 β a))
\end{verbatim}
}

\section{Proof strategy}
On $(1,L]$ the integrand $H(x)/x$ is rewritten pointwise as
$A\log x/x+C_0x^{-1}+\theta(x)/x$; the first two integrate to $\tfrac A2\log^2L+C_0\log L$ by the
FTC lemmas of unit~38, and the third is $\int_1^\infty\theta/x-\int_L^\infty\theta/x$ by
\texttt{setIntegral\_union}. Combined with the split at $x=1$ this is the exact identity. The
integrability of $\theta(x)/x$ on $[1,\infty)$ uses the domination $0\le\theta/x\le Mx^{-2}$ from
unit~41 and \texttt{integrableOn\_Ioi\_rpow\_of\_lt}; measurability comes from continuity on
$(1,\infty)$, where $H$ agrees with its interval-integral primitive and $\log$ is continuous. The
$Z_3$ statements then follow from $Z_3=n^{-1/2}H_2(\sqrt nb^3)$ and
$\log(\sqrt nb^3)=\tfrac12\log n+3\log b$ exactly as at $d=2$; the polynomial coefficient is kept in
nested form and expanded by a separate \texttt{ring} lemma so the \texttt{Tendsto} statement can name
the constant. The identity and remainder sign were checked numerically first (residual $0$ at $L\ge5$
to double precision).

\section{Roadblocks and resolutions}
\begin{itemize}
\item \texttt{Integrable.add} returns a \texttt{Pi.add}; \texttt{rw [integral\_add …]} then fails to
match the $\lambda$-integrand. Type-ascribe the combined integrability witness as a single lambda
(the documented workaround) before rewriting.
\item \texttt{positivity} cannot see $0\le M$ for the opaque \texttt{quadMoment}; supply
\texttt{quadMoment\_nonneg} through \texttt{mul\_nonneg}/\texttt{div\_nonneg}.
\item \texttt{field\_simp} on the pointwise decomposition of $H(x)/x$ leaves a commutativity goal;
follow with \texttt{ring}.
\end{itemize}

\section{What was Mathlib, what was new}
Mathlib: \texttt{ContinuousOn.aestronglyMeasurable}, \texttt{integrableOn\_Ioi\_rpow\_of\_lt},
\texttt{integral\_Ioi\_rpow\_of\_lt}, \texttt{setIntegral\_union}, \texttt{squeeze\_zero'}. New: the
remainder $\theta$ and constant $C_1$ of $H_2$, and the complete $O(n^{-1})$-accurate expansion of the
$d=3$ log-squared chart integral.

\section{Lessons}
The one-step transfer ``polynomial in $\log$ plus $O(1/x)$'' $\mapsto$ same shape one degree higher
is now visible as a reusable pattern; a general-$m$ version for \texttt{iterLogPrim} would make the
constants below $(\log n)^{d-1}$ available for every $d$ at once, mirroring how unit~39 generalised
units~35 and~38.

\section{Outcome}
The $d=2$ and $d=3$ coinciding-exponent integrals now carry complete polynomial-in-$\log n$
expansions with explicit constants and $O(n^{-1})$ remainders. Next candidates: the general-$m$
transfer lemma for \texttt{iterLogPrim}, or the weighted primitives $F_\gamma$ for general $(k,h)$.
\end{document}
